\section{Domaines et conditions des réalisations}
\label{v:integracion}
L’exclusion algébrique, son transport par raffinement et ses applications
statistiques sont des résultats liés, dont les domaines sont distincts.
L’identité générale d’annihilation
\eqref{v:fock:eq:transporte-aniquilacion} explicite l’adjoint du transport ;
sa spécialisation isométrique conserve la condition nécessaire portant sur
l’application entre les espaces à une particule. Les occupations
d’équilibre utilisent, en outre, une graduation et un Hamiltonien compatibles
avec la représentation de la trajectoire.

La réalisation dynamique finie est développée dans la
section~\ref{sec:hmtf2-cristal-temporal}. L’opérateur d’aire, le
Hamiltonien modulaire et la purification sectorielle sont développés dans la
section~\ref{v:ant:sec:area} et ses prolongements. Les sections
\ref{v:ant:sec:action} et \ref{v:ant:sec:ii-kb-aut-desarrollo} contiennent
les constructions de l’action et du couple Boltzmann--température qu’utilise
la réalisation radiative. La caractérisation variationnelle de la
mesure des histoires et la convergence des sommes spectrales relient,
respectivement, l’incidence à l’entropie et les modes discrets au rayonnement.
La section radiative
présente ses hypothèses de cavité,
de dispersion, de polarisation et d’équilibre, ainsi que les preuves spectrales
qu’elle utilise. La section~\ref{v:sec:enlace-velocidad-constitutiva}
réunit en outre l’identité des réalisations cinématique et constitutive
de la vitesse, sa normalisation dimensionnelle et sa conservation sous
amplification de l’espace à deux feuillets.

Chaque composition conserve les hypothèses mathématiques et la réalisation
physique qu’elle déclare. En particulier, les modes harmoniques du complexe
cellulaire, les connexions modulo jauge et les modes radiatifs conservent
les domaines distingués dans la
section~\ref{v:rec:completaciones:section:08} ; la réalisation radiative
utilise explicitement sa dispersion et ses polarisations.

La sous-section~\ref{v:subsec:limite-termodinamico} démontre, pour les trois
observables \(F_N,F_E,F_Z\), le passage des sommes périodiques tridimensionnelles
aux intégrales radiales lorsque \(L\to\infty\). Son domaine exige \(a,b>0\),
maintient fixé le mode constant de fond et contrôle uniformément l’origine
et la queue spectrale. Le résultat fournit les densités de nombre,
d’énergie d’excitation et de logarithme de la fonction de partition des modes
transversaux ; il est distinct du bloc fini de capacité spectrale.

La section~\ref{v:sec:capacidad-espectral-termica} a pour domaine le
support gradué fini, les arbres microscopique et agrégé, le
relèvement réversible de mémoire, l’horloge de capacité et les systèmes de coordonnées
thermométriques qui y sont déclarés. Ses identités spectrales, de transport et
variationnelles sont prouvées dans ces domaines ; la reconnaissance physique
intervient ensuite. Le système de coordonnées dimensionnelles typé conserve la provenance de
l’intensité marquée, du système de coordonnées énergétiques et de la température physique.

La section~\ref{v:n19:sec:incertidumbre-fourier} travaille sur la
réalisation finie \(\ell^2((\mathbb Z/27\mathbb Z)^2)\) et sa
transformée de Fourier unitaire. Ses entropies sont celles des distributions dans deux bases
conjuguées ; les lecteurs dodécaphasiques possèdent le domaine différent
\(\{-1,0,+1\}^{12}\) et apportent un résultat complémentaire de mémoire
ordonnée. La section~\ref{sec:f051-prime-vacancy-observer} établit
l’identité modale exacte pour chaque \(X\) fini ; son témoin historique atteint
\(10^8\) et la reproduction indépendante documentée atteint \(10^6\).
La promotion à la borne critique requiert des estimations uniformes sur les modes.
